Was ist ein Versicherungsvertrag im Sinne des IFRS 17?

Ein Versicherungsvertrag ist “ein Vertrag, nach dem eine Partei (der Versicherer) ein signifikantes Versicherungsrisiko
von einer anderen Partei (dem Versicherungsnehmer) übernimmt.”
Im Rahmen des Vertrages wird vereinbart, dass der Versicherer dem Versicherungsnehmer eine Entschädigung
zahlt, “wenn ein festgelegtes ungewisses künftiges Ereignis (das versicherte Ereignis) den Versicherungsnehmer
nachteilig betrifft.“


Welche Aspekte fallen nicht unter IFRS 17 (Anwendungsausschlüsse)?

Bestimmte Garantien
Vermögenswerte und Verbindlichkeiten des AG aus Versorgungsplänen
Finanzgarantien unter bestimmten Voraussetzungen
Bestimmte Restwertgarantien
Vertragliche Rechte oder Pflichten, die von der künftigen Verwendung oder dem Recht zur Nutzung eines nicht-
finanziellen Postens abhängen
Im Rahmen eines Unternehmenszusammenschlusses zu zahlende oder ausstehende bedingte Gegenleistungen
Kreditkarten mit Versicherungsschutz
Erstversicherungsverträge beim VN

Welche Komponenten müssen von einem Versicherungsvertrag im Sinne des IFRS 17 separiert werden?

VersicherungsKomponente IFRS-17
Nicht-abgrenzbare Investment-komponente IFRS 15
Abgrenzbare Investmentkomponente IFRS 9
Gesonderte Dienstleistungen IFRS 17
Nicht eng verbundenes eingebettetes Derivat IFRS 17, kein GuV- Ausweis der Ein- und Auszahlungen

Eine Investmentkomponente ist eindeutig abgrenzbar, wenn sie nicht in hohem Maße verbunden ist und ein Vertrag
mit vergleichbaren Konditionen einzeln verkauft werden könnte


Welche vier Grundbestandteile hat das allgemeine Bewertungsmodell?

Das allgemeine Bewertungsmodell stützt sich auf die Methode "Building Block Approach"
was Bestandteil des General measurement models ist.
(die Vorlesung nennt es "Erstmaliger Ansatz", bisher noch nicht woanders gefunden).
Er besteht aus
Block 1: Die Summe aller zukünftigen Zahlungsströme.
Block 2: Der diskontierte Wert dieser Zahlungsströme.
Block 3: Risikoanpassung (Risk adjustment), spiegelt die Einschätzung des Unternehmens zu der wirtschaftlichen Last wider, 
         die ein spezifisches Risiko in sich trägt
Block 4: Vertragliche Servicemarge/ Contractual service margin (CSM), Entspricht den unrealisierten erwarteten Gewinnen (in der Zukunft) – 
         dadurch wird das Imparitätsprinzip gewahrt und Gewinne erst verbucht, wenn sie wirklich gezahlt werden. 

1+2+3 = Erfüllungswert = FCF fullfilment cash flows

INSERT IMAGE 26

Wie ist das Aggregationsniveau bei der Bewertung?

Bewertungseinheit der CSM ist die Gruppe von Versicherungsverträgen (GVV) welche nach min. 3
aus dem PVV gefiltert wurden.

Das bedeutet von einem Portfolio von Versicherungsverträgen (PVV) (alle welche die Versicherungsparte hat),
wird eine zuerst eine Oberart von Verträgen gewählt, z.B. Risikolebensversicherungsverträge.
Dann wird diese in Kohorten aufgeteilt (Erstellungsjahr der Verträge).
Zum schluss werden sie in Defizitäre, Sichere und Sonstige Verträge aufgeteilt was dann die GVV sind.

Innerhalb einer GVV sollen die Zeichnungsdaten
der Verträge maximal 12 Monate
auseinanderfallen.

Geben Sie ein Beispiel zur Kalkulation des "Building Block Approaches"


INSERT IMAGE 26 (IFRS17_CSMWorkingParty_SessionalPaper)

Wir sehen dass wenn der Versicherungsvertrag defizitär ist dann fällt die CMS auf 0.
(Beantwortet Welche Besonderheit gibt es bei verlustträchtigen Verträgen?)


Wie erfolgt die Folgebewertung grundsätzlich?

Da die Gewinne welche in die CSM inkludiert wurden über den Zeitverlauf in realisierte Gewinne umgewandelt
werden müssen, ist eine Folgebewertung unabdingbar. Zusätzlich sind neue Zinssätze, höhere Schadensquoten etc.
weitere Punkte für eine Folgebewertung.

Die Höhe des in der GuV anzusetzenden Betrages aus der Auflösung der CSM bestimmt sich aus der Anzahl der in der Periode
und in zukünftigen Perioden zu leistenden Deckungseinheiten/Serviceeinheiten.

Die Menge der aktuellen Deckungseinheiten bestimmt im Verhältnis zu den aktuell zu leistenden Deckungseinheiten in der
Periode sowie derer über die verbleibende Laufzeit den aufzulösenden Anteil der CSM

INSERT IMAGE 13 (IFRS17_CSMWorkingParty_SessionalPaper)


Als weiteres Beispiel, am Ende des Geschäftsjahres wird vermutet, dass die Schadenzahlungen für das kommende Jahr 550 anstatt 200 betragen werden. Die
Bewegungen der Verbindlichkeiten sind wie folgt:

INSERT IMAGE 50



Welche Bestandteile hat der Buchwert einer Gruppe von Versicherungsverträgen und 
auf welcher Ebene werde diese in der Bilanz ausgewiesen?

Buchwert = LRC + LIC

LRC = Liability for remaining coverage (vt. Deckungsrückstellung = CMS + Erfüllungswert)
LIC = Liability for incurred claims (Rückstellung für noch nicht abgewickelte Versicherungsfälle)

Portfolien von Versicherungsverträgen, die Vermögenswerte sind
Portfolien von Versicherungsverträgen, die Verbindlichkeiten sind
Rückversicherungsverträge werden davon gesondert ausgewiesen, ebenso danach getrennt, ob es sich um einen
Vermögensgestand oder eine Verbindlichkeit handelt.

Welche Komponenten sind in der Gewinn- und Verlustrechnung zu berücksichtigen?

In der GuV wird das Versicherungstechnische Ergebnis/Insurance service result ausgewiesen.
Es besteht aus
Versicherungsertrag/Insurance revenue 
(z.B. Auflösung der Deckungsrückstellung, d.h. erwartete Cashflows, Risk Adjustment, CSM)
Versicherungsaufwand/Insurance service expenses 
(z.B. eingetretene Leistungen und Kosten, Veränderungen von Leistungen und Risk Adjustment)


Welche Ziele verfolgen die Anhangangaben?

den Nutzer dabei unterstützen, die Auswirkungen der Verträge beurteilen zu
können, in Bezug auf Finanzlage, Finanzielle Leistung und Zahlungsströme

Notwendige Angaben sind:
Überleitungsrechnungen zur Darstellung, wie sich die Buchwerte über den Berichtszeitraum infolge von Zahlungsströmen,
Erträgen und Aufwendungen, die in der Periode entstanden sind, verändert haben.

Erläuterung, wann das Unternehmen damit rechnet, die zum Berichtszeitpunkt verbleibende vertragliche Servicemarge
erfolgswirksam in der GuV zu erfassen.

Ermessensentscheidungen in der Bewertung und Informationen über die Annahmen/den Input

Angaben zu Risiken (vor allem auf Versicherungs- und Finanzrisiken ).


Welche Besonderheit gibt es bei der Bewertung mit der Methode des Prämienallokationsansatz und welche Anforderungen gibt
es an die Verträge?

Die Methode des Prämienallokationsansatz (Premium Allocation Approach, PAA) ist ein optionales, vereinfachtes
Bewertungsmodell für die Bewertung der Deckungsrückstellung (LRC).

Im Gegensatz zum Building Block Approach wird die LRC nicht über den present value der FCF + CMS bestimmt, 
sondern schlicht als ein fixer Anteil (bisher vergangene Zeit/gesamte Vertragslaufzeit) der Nettopremiumzahlung
minus Abschlusskosten plus Verlustkomponente.

Hinsichtlich der LIC werden weiterhin Risk adjustment sowie present value der FCF verwendet.

Ob die PAA genutzt werden darf leitet sich aus folgendem Diagram ab

INSERT IMAGE 69

Eine vernünftige Annäherung an das GMM ist nicht gegeben, wenn das Unternehmen eine signifikante Variabilität der
Erfüllungswerte während der Versicherungsdauer erwartet. Z.b.
- Erfüllungswerte während der Versicherungsdauer erwartet
 dem Umfang der Zahlungsströme aus eingebetteten Derivaten im Vertrag.
- der Länge der Versicherungsdauer.
Beispiele für Sparten mit PAA-Anwendung sind: Kfz, private Haftpflicht, Hausrat, Gebäude.


Welche Besonderheit gibt es bei der Bewertung der direkt beteiligten Verträge und welche Anforderungen gibt es an
die Verträge?

Für die Bilanzierungsmodelle zur Bewertung von Lebensversicherungsverträgen sind zwei Methoden vorgesehen: 
das GMM und der Variable Fee Approach (VFA). 
Letzterer findet Anwendung bei der Bilanzierung von Verträgen mit wiederbewertbaren Leistungen auf Basis der Gewinnbeteiligung. 
Solche Verträge werden auch direkt beteiligte Verträge genannt.
Der VFA ist eine Variante des GM und wird angewendet, um Provisionen für die Verwaltung von Vermögenswerten (Verwaltungsgebühren) zu berücksichtigen, 
die an das Versicherungsunternehmen gezahlt werden.
Das Verfahren reduziert die Volatilität der Nettoresultate.

Direkt beteiligte Verträge sind erkennbar daran, dass der Versicherungsnehmer einen
Anteil eines klar identifizierten Bezugswertes erhält.
Der identifizierte Bezugswert muss nicht vom VU gehalten werden. Beispiele sind die Leben- und Krankenversicherung.



Welche Besonderheit gibt es bei der Bewertung von passiver Rückversicherung?

Rückversicherungen werden getrennt ausgewiesen in der GuV, aber dennoch nach der PPA bzw. GMM 
wie das zedierten Geschäft bewertet.

Alle anderen Zahlungsströme, die nicht unmittelbar eine Zession
darstellen, werden eigenständig angesetzt und bewertet.
Der Erfüllungswert beinhaltet das Nichterfüllungsrisiko des
Rückversicherers, Änderungen dieses Risikos sind nicht den
zukünftigen Leistungen zuzuordnen sondern direkt in der GUV.

Die Risikoanpassung bildet die Höhe des kontinuierlich
bewerteten Risikos ab, das von dem Zedenten auf den
Rückversicherer übertragen wird – dementsprechend erhöht
die Risikoanpassung den aktivischen Wert der
Rückversicherung

Die CSM repräsentiert die Netto-Kosten oder den Netto-Ertrag aus dem Erwerb einer passiven Rückversicherung.
Die Netto-Kosten einer Rückversicherungsnahme für Ereignisse, die sich vor dem Abschluss ereignet haben, werden unmittelbar als
Aufwand erfasst (retroaktive Rückversicherung).



Ab wann gilt die Anwendung von IFRS 17?

Grundsätzlich zum Übergangsstichtag. Dabei wird jede GVV so bewertet als hätte IFRS 17 schon immer gegolten und entfernt
Posten die nicht eingerechnet werden können.


Die vollständige rückwirkende Anwendung von IFRS 17 ist möglich, wenn vernünftige und belastbare Informationen, 
die ohne unangemessene Kosten oder Mühen verfügbar sind. Es dies nicht der Fall dürfen in soweit Modifikationen genutzt werden,
als das ein Ergebnis erzielt werden, das möglichst nah an dem von IFRS 17 liegt.

Diese Entscheidung kann pro GVV getroffen werden.

Die Modifikationen konzentrieren sich auf die Schätzung der CSM bzw. der Verlustkomponente um die LRC zum Übergangszeitpunkt zu bestimmen.
Modifikationen sind vorgesehen zur Bestimmung von
• der GVV und zur Klassifizierung von Verträgen beim erstmaligen Ansatz
• der Bestandteile des allgemeinen Bewertungsmodells, einschließlich der Verlustkomponente
• Beträge die durch Aufteilung auf GuV und OCI enstehen.

Alternativ wird der Zeitwert-Ansatz genutzt
dieser hat seinen Ursprung in dem Fakt, dass die Höhe der CSM oder der Verlustkomponente am Übergangsstichtag basiert auf der zu diesem Zeitpunkt
bestehenden Differenz zwischen dem beizulegenden Zeitwert und dem Erfüllungswert der GVV


Welche Methoden gibt es in der Ermittlung der Eröffnungsbilanz?
Vollständig rückwirkender Ansatz
Modifizierter rückwirkender Ansatz,
Zeitwert-Ansatz

Es muss erkennbar sein, welche Methode zur Bestimmung der Übergangswerte verwendet wurde und wie
Beurteilungsspielräume genutzt wurden.
